\section{项目经历}
列举部分早期项目研究经历如下。 欢迎到我的 \href{https://github.com/RWLinno}{Github} 关注最新项目动态!

\cvitem{大创省级}{\textbf{小麦价格组合预测模型}(2020-2023) - 基于ARIMA、LSTM、XGBoost构建多元时序组合预测模型。}
\cvitem{挑战杯校级}{\textbf{基于ViT的医学影像辅助诊断系统} - 应用Vision Transformer处理CT/MRI/X光片，GitHub 10+ Stars。}
\cvitem{挑战杯校级}{\textbf{退行性骨关节病AI辅助诊断} - 基于中医数据集的多标签分类任务，结合特征工程与深度学习。}
\cvitem{}{\textbf{IVDPS:集成视觉数据处理系统} - 支持数字图像处理、情感识别等多功能的可视化系统。}
\cvitem{OSM-Artist}{\textbf{地理信息可视化工具} - 基于osmnx和OpenStreetMap的城市路网快速绘图工具。}
\cvitem{SojumpHelper}{\textbf{问卷互填助手} - 问卷星互填社区快速得分工具，帮助快速收集样本。}
\cvitem{}{\textbf{图书管理及学生选课系统} - 基于.NET/C\#/MySQL/WinForms的GUI程序设计，MySQL数据库增删改查设计与开发。}
% 下面这些不需要
% \cvitem{论文代码}{\textbf{ST-LoRA / STUM} - 时空预测论文的官方PyTorch实现，合计GitHub 13+ Stars。}
% \cvitem{学术资源}{\textbf{ST-Papers-with-Slides} - 收集整理时空数据挖掘顶会论文Slides，便于学术交流学习。}
\cvitem{DevOps-MAS}{\textbf{多智能体运维系统} - 基于VLM+RAG+MCP的智能运维多智能体系统，集成专业智能体协同、三层记忆架构、知识库检索增强与SFT数据流水线；支持OpenAI/Ollama/vLLM多后端，提供Streamlit GUI与Docker一键部署。}

\cvitem{ST-LoRA}{\textbf{面向时空预测的参数高效微调框架} - 通过节点自适应低秩预测器实现<1\%额外参数带来~7\%性能提升；支持STGNNs即插即用，提供完整训练/推理/可视化流水线。}

%\cvitem{STUM}{\textbf{Spatio-Temporal Unitized Model for Traffic Flow} (IEEE T-ITS, JCR Q1) - 时空交互统一建模框架，提出自适应时空单元聚合模块（ASTUC）与多粒度局部感受野（MLRF），在PEMS系列数据集上实现SOTA；代码开源含数据预处理、基准对比与教程。\textit{PyTorch, Jupyter, 5⭐}}
\cvitem{Agentic VQA Cleaner}{\textbf{具身VQA数据蒸馏/清洗/合成} - 多模态数据合成与校验工具，支持图像质量/问题合理性/答案准确性三维评估，兼容2D/3D场景图推理、启发式规则与并发处理。}

%\cvitem{ST-Papers-with-Slides}{\textbf{Spatio-Temporal Papers Collection with Slides} - 持续更新的时空数据挖掘顶会论文库，整理AAAI/KDD/ICML等会议论文+官方Slides+复现笔记，助力领域快速入门与学术交流。\textit{Markdown, 3⭐}}

% 早期教学/工具类项目归档（可按需展开）
% \cvitem{教育工具}{SojumpHelper问卷助手 / OSM-Artist地理可视化工具 / 图书选课系统(.NET) 等，详见个人主页}
